\documentclass[../../../include/open-logic-section]{subfiles}

\begin{document}

\olfileid{sth}{choice}{hartogs}
\olsection{مقایسه‌پذیری و لم هارتوگس}

این جنبهٔ مثبت ماجرا بود. اکنون جنبهٔ منفی آن را ببینیم. بدون اصل
انتخاب، اوضاع \emph{آشفته} می‌شود. برای روشن‌شدن علت، نتیجهٔ زیبای
زیر را از \cite{Hartogs1915} در نظر بگیرید:

\begin{lem}[\emph{در $\ZF$}]\ollabel{HartogsLemma}
برای هر مجموعهٔ $A$، عدد ترتیبیِ $\alpha$ای وجود دارد به‌طوری‌که
$\cardnless{\alpha}{A}$.
\end{lem}

\begin{proof}
اگر $B \subseteq A$ و $R \subseteq B^2$، آن‌گاه بنا بر
\olref[ord-arithmetic][using-addition]{rankcomputation} داریم
$\tuple{B, R} \subseteq V_{\setrank{A}+4}$. پس با استفاده از
جداسازی، در نظر بگیرید:
\[
	C = \Setabs{\tuple{B, R} \in V_{\setrank{A}+5}}{B\subseteq A
	\text{ و $\tuple{B, R}$ یک خوش‌ترتیب است}}
\]
با استفاده از جایگزینی و
\olref[ordinals][ordtype]{thmOrdinalRepresentation} مجموعهٔ زیر را
تشکیل دهید:
\[
	\alpha = \Setabs{\ordtype{B, R}}{\tuple{B, R} \in C}.
\]
بنا بر \olref[ordinals][basic]{corordtransitiveord}، $\alpha$ عددی
ترتیبی است، زیرا مجموعه‌ای متعدی از اعداد ترتیبی است. به‌راستی، اگر
$\gamma \in \beta \in \alpha$، آن‌گاه برای یک
$B \subseteq A$ داریم $\beta = \ordtype{B, R}$؛ سپس بنا بر
\olref[ordinals][iso]{wellordinitialsegment}، برای یک $b \in B$
داریم $\gamma = \ordtype{B_b, R_b}$، و بنابراین
$\gamma \in \alpha$.

برای برهان خلف، فرض کنید !!a{injection} مانند
$f \colon \alpha \to A$ وجود داشته باشد. تعریف کنید:
\begin{align*}
	B &= \ran{f}\\
	R &= \Setabs{\tuple{f(\alpha), f(\beta)} \in A \times A}{\alpha \in \beta}.
\end{align*}
آشکار است که $\alpha = \ordtype{B, R}$ و
$\tuple{B, R} \in C$. پس $\alpha \in \alpha$، که تناقض است.
\end{proof}

از این لم نتیجه‌ای ژرف به دست می‌آید:

\begin{thm}[\emph{در $\ZF$}]
ادعاهای زیر هم‌ارزند:
\begin{enumerate}
	\item\ollabel{equivwo} اصل موضوع خوش‌ترتیبی؛
	\item\ollabel{equivcompare} برای هر دو مجموعهٔ $A$ و $B$، یا
	$\cardle{A}{B}$ یا $\cardle{B}{A}$.
\end{enumerate}
\end{thm}

\begin{proof}
\emph{\olref{equivwo} $\Rightarrow$ \olref{equivcompare}.}
$A$ و $B$ را ثابت کنید. با به‌کارگیری \olref{equivwo}،
خوش‌ترتیب‌هایی مانند $\tuple{A, R}$ و $\tuple{B, S}$ وجود دارند.
با به‌کارگیری
\olref[ordinals][ordtype]{thmOrdinalRepresentation}، بگذارید
$f \colon \alpha \to \tuple{A, R}$ و
$g \colon \beta \to \tuple{B, S}$ یکریختی باشند. بنا بر
\olref[sth][ordinals][basic]{ordinalsaresubsets}، یا
$\alpha \subseteq \beta$ یا $\beta \subseteq \alpha$. اگر
$\alpha \subseteq \beta$، آن‌گاه
$\comp{f^{-1}}{g} \colon A \to B$ یک !!a{injection} است و ازاین‌رو
$\cardle{A}{B}$؛ به همین ترتیب، اگر
$\beta \subseteq \alpha$، آن‌گاه $\cardle{B}{A}$.

\emph{\olref{equivcompare} $\Rightarrow$ \olref{equivwo}.}
$A$ را ثابت کنید؛ بنا بر \olref{HartogsLemma} عدد ترتیبیِ
$\beta$ای وجود دارد که $\cardnless{\beta}{A}$. با به‌کارگیری
\olref{equivcompare}، داریم $\cardle{A}{\beta}$. پس
!!{injection} مانند $f \colon A \to \beta$ وجود دارد، و می‌توانیم
با تعریف ترتیب
$\Setabs{\tuple{a, b} \in A \times A}{f(a) \in f(b)}$، از این تزریق
برای خوش‌ترتیب‌کردن اعضای $A$ استفاده کنیم.
\end{proof}
\noindent
در نتیجه‌ای بی‌درنگ، اگر خوش‌ترتیبی شکست بخورد، برخی مجموعه‌ها از
نظر اندازه \emph{به‌معنای واقعی کلمه مقایسه‌ناپذیر} خواهند بود.
بنابراین، در صورت شکست خوش‌ترتیبی، حساب کاردینالیِ ترامتناهی آشفته
می‌شود. برای نمونه، باید از این اندیشه دست بکشیم که اگر $A$ و $B$
نامتناهی باشند، آن‌گاه $\cardeq{A \disjointsum B}{M}$ و
$\cardeq{A \times B}{M}$، که در آن $M$ بزرگ‌ترِ $A$ و $B$ است
(نگاه کنید به
\olref[card-arithmetic][simp]{cardplustimesmax}). مشکل ساده است:
اگر نتوانیم اندازهٔ $A$ و $B$ را \emph{مقایسه} کنیم، پرسیدن اینکه
کدام‌یک بزرگ‌تر است بی‌معناست.

\end{document}
